% Generated from locale/tr/content/lambda-calculus/introduction/overview.tex; do not hand-edit.


\olfileid[tr]{lam}{int}{ovr}
\olsection{Genel Bakış}

Lambda kalkülüsü ilk olarak Alonzo Church tarafından 1930'ların başında
yapıcı mantığın temeli olarak tasarlanmıştı; hesaplanabilir fonksiyonların
bir modeli olarak \emph{tasarlanmamıştı}. Ancak kısa süre sonra Turing
hesaplanabilir fonksiyonları ve kısmi özyinelemeli fonksiyonlar gibi diğer
hesaplanabilirlik tanımlarına eşdeğer olduğu gösterildi. Bunun başlangıçta
biraz şaşırtıcı olması, nitelendirmeyi daha da ilginç kılar.

Lambda gösterimi, bir fonksiyona ad tanımlamak yerine, fonksiyonu tanımlayan
simgesel bir ifadeyle doğrudan gönderme yapmanın elverişli bir yoludur.
``$f$, $f(x)=x+3$ ile tanımlanan fonksiyon olsun'' demek yerine, ``$f$,
$\lambd[x][(x+3)]$ fonksiyonu olsun'' denebilir. Başka bir deyişle
$\lambd[x][(x+3)]$, argümanına üç ekleyen fonksiyonun yalnızca bir
\emph{adıdır}. Bu ifadede $x$ bir kukla değişken ya da yer tutucudur: aynı
fonksiyon $\lambd[y][(y+3)]$ ile de gösterilebilir. Gösterim başka
parametreler varken de çalışır. Örneğin $g(x,y)$ iki değişkenli bir fonksiyon
ve $k$ bir doğal sayı olsun. O zaman $\lambd[x][g(x,k)]$, her $x$'i
$g(x,k)$'ye götüren fonksiyondur.

Bir fonksiyonun simgesel bir ifadeden bu biçimde tanımlanmasına
\emph{lambda soyutlaması} denir. Lambda soyutlamasının öbür yüzü
\emph{uygulamadır}: bir $f$ fonksiyonu (diyelim doğal sayılar üzerinde
tanımlı) verildiğinde, fonksiyon $2$ gibi herhangi bir değere
\emph{uygulanabilir}. Alışılmış gösterimde sonuç elbette $f(2)$ ile yazılır.

Lambda soyutlamasını uygulamayla birleştirince ne olur? Ortaya çıkan ifade,
uygulanılan ifadeyi soyutlanan değişkenin yerine ``koyarak'' sadeleştirilebilir.
Örneğin
\[
(\lambd[x][(x + 3)])(2)
\]
ifadesi $2+3$'e sadeleştirilebilir.

Buraya kadar alışılmış kavramlar için yeni gösterimler sunmaktan başka bir
şey yapmadık. Fakat lambda kalkülüsü, küme kuramsal bakış açısından daha
köklü bir kopuşu temsil eder. Bu çerçevede:
\begin{enumerate}
\item Her şey bir fonksiyon gösterir.
\item Fonksiyonlar lambda soyutlamasıyla tanımlanabilir.
\item Her şey başka her şeye uygulanabilir.
\end{enumerate}
Örneğin $F$, lambda kalkülüsünde bir terimse $F(F)$ ifadesinin her zaman
anlamlı olduğu varsayılır. Bu özgür çerçeveye \emph{türsüz} lambda kalkülüsü
denir; burada ``türsüz'', ``neyin neye uygulanabileceği üzerinde hiçbir
kısıtlama yok'' demektir.

\begin{digress}
Türsüz sürümün önemli bir değişkesi olan \emph{türlü} lambda kalkülüsü de
vardır. Türlü lambda kalkülüsü birçok bakımdan türsüz olana benzese de klasik
küme kuramsal çerçeveyle bağdaştırılması çok daha kolaydır ve oldukça farklı
bazı özelliklere sahiptir.

Lambda kalkülüsü araştırmaları kuramsal bilgisayar biliminde ve programlama
dillerinin tasarımında merkezi önem kazanmıştır. John McCarthy'nin 1950'lerde
tasarladığı LISP, bu düşüncelerden etkilenen ilk dillerden biridir.
\end{digress}

